\documentclass[../main]{subfiles}
\begin{document}
\chapter*{Planificación 2026}
\info{Planificación de la asignatura: \emph{Dibujo Técnico Asistido}}
{
Escuela 4-063 "Luis Federico Leloir"

Espacio Curricular: Dibujo Técnico Asistido.

Docente: Ferreyra Gustavo David.

Curso: $3^{ro} 2^{da}$ 
}
\section{Acuerdos institucionales generales}

En primer lugar, se expresa que se dará cumplimiento las distintas pautas y acuerdos establecidos de en forma institucional tales como:
\subsection*{Ejes transversales}
\begin{itemize}
    \item  Manejar e interpretar documentación técnica, protocolos y/o manuales específicos pertinentes a la actividad electrónica.
    \item Elaborar conclusiones a partir de la reflexión critica y la evaluación vinculando los saberes teóricos con los prácticos.
\end{itemize}
\subsection{Competencia comunicativa}
\subsubsection{Fluidez lectora}
\begin{itemize}
    \item Realizar una lectura en voz alta, 10 minutos por semana en cada espacio curricular.
    \item Proponer a los estudiantes grabarse en audios, de tal manera que autoregulen sus procesos de lectura.
    \item Enseñar el sentido de los distintos signos de puntuación.
\end{itemize}
\subsubsection{Comprensión lectora}
\begin{itemize}
    \item Incluir conscientemente la lectura de los paratextos (por eje dibujos, infografías, mapas,etc)
    \item Trabajar los títulos, substitulos y fuentes.
    \item Corregir la ortografía y signos de puntuación.
    \item Ejercitar la práctica de extraer ideas principales contenidas en los párrafos.
    \item Traer a colación ejemplos de la vida contidiana y de los intereses de los adolescentes, para establecer relaciones con lo emocional que asegura la perdurabilidad de los aprendizajes.
    \item Plantear exposiciones orales de modo más frecuente.
    \item Elaborar glosario en todas las asignaturas.
    \item Trabajar modelos de informes, consensuados con los profesores de Lengua.
    \item Realizar preguntas inferenciales.
    \item "Receta infalible" leer, identificar, explicar los comprendido, transferir y volver a repetir.
\end{itemize}
\subsection{Trabajar en A.B.P}
\begin{itemize}
   \item En lo posible realizar articulaciones con otras instituciones, empresas, actores de la salud y sociales.
 \item Fortalecer el trabajo por CAPACIDADES, entendiendo que son transversales y se logran con un trabajo en equipo. Potenciar el trabajo institucional en proyectos significativos para los adolescentes, asi como pertinentes a su perfil profesional.
 \item fomentar las evaluaciones por rúbricas y el trabajo colaborativo entre docentes. Incorporar las nuevas tecnologías como un recurso para la enseñanza y el aprendizaje. trabajar con simulaciones, programas existentes para los distintos espacios, tutoriales (del profesor del espacio o los existentes en redes.
\end{itemize}
\subsection{Mantener una fluida comunicación}
Entre todos los sectores de la escuela.  La eficienca en la comunicación impactó en los buenos rendimientos de los estudiantes y en aspectos preventivos de los confictos.
Profundizar los trabajos en equipo, el diálogo y los acuerdos que dan estabilidad y previsibilidad a la tarea docente.
\section{Acuerdos específicos del área}
\subsection{Acuerdos 2025: Proyecto de integración del pensamiento computacional, breve descripción:}
Este es un proyecto que pretende modificar de manera sustancial la forma en que se desarrollan las prácticas en los talleres de electrónica. La electrónica está cada vez más integrada con la programación, y hoy no se concibe una electrónica que no tenga algún grado de programación. Es por esto que respetando los contenidos del diseño curricular provincial, pero atendiendo a esta realidad que atraviesa nuestra especialidad es que se acordó el área incluir en todos los años algo de programación. La organización es.
3 año: Programación de la impresora 3D. Se enseñará a trabajar en el manejo, calibración, control y realización de piezas en la impresora 3D Para el espacio \emph{introducción a las Prácticas Productivas en Electrónica} se desarrollará una parte de Electrónica y Audio, y con la adecuación de los horarios de dos de los maestros de enseñanza práctica, se implementará el área de programación de impresoras 3D. En este año también se quiere vincular el trabajo de impresora 3D con Diseño Asistido por Computadora.
\subsection{Reajuste de Prácticas Profesionalizantes.}
Hasta el año pasado se trabajaba en las prácticas profesionalizantes un proyecto con su producción en quinto año, y otro trabajo con su producido en sexto año. A partir de esto se acordó que existiese un proyecto en común que comience en quinto año y termine en sexto año. basado en la automatización de procesos, equipos o dispositivos, haciendo uso de todo lo visto en la carrera y con centro en PLC y arduino. Para esto se ha organizado que en las prácticas profesionalizantes de quinto se desarrolle las siguientes actividades.

\section{Acuerdos del área}
\subsection{Acuerdo por capacidades}
Comprensión y producción de textos escritos y exposiciones orales inherentes a temas del área.
Abordaje y resolución de situaciones problematicas en forma analítica, gráfica y práctica.
Trabajo en colaboración para aprender a relacionarse e interacutar.
\subsection{Acuerdo de actividades y procedimientos}
Exigirá la realización y presentación de una carpeta técnica individual. La misma debe tener una portada formal, índice, planificación o programa de estudio, 
% ####################### Bloque 1- #######################################
\begin{table}[h]
    \centering
    \begin{tabular}{|c|c|} \hline
         \multicolumn{2}{|c|}{Contenidos} \\ \hline
         Conceptuales & Procedimentales  \\ \hline \hline
        \begin{minipage} [t] {0.4\textwidth} 
            \begin{itemize}
                \item Letras. Formato de Láminas.
                \item Líneas, Aplicaciones. 
                \item Interrupción de cuerpos.
                \item Líneas de rotura,aplicaciones.
                \item Rayados y colores convencionales.
                \item Modos de plegar planos.
                \item Rotulación de planos.
            \end{itemize}
        \end{minipage} &  \begin{minipage} [t] {0.4\textwidth} 
            \begin{itemize}
                \item Representar cuerpo aplicando diferentes tipos de línea.
                \item Representación de dibujos en láminas con formato y rótulo adecuado.
                \item Plegado de planos.
            \end{itemize}
        \end{minipage} \\ \hline
       \end{tabular}
    \caption{Bloque 1}
\end{table}


\begin{table}[h]
    \centering
    \begin{tabular}{|c|c|} \hline
         \multicolumn{2}{|c|}{Contenidos} \\ \hline
         Conceptuales & Procedimentales  \\ \hline \hline
        \begin{minipage} [t] {0.4\textwidth} 
            \begin{itemize}
                \item Métodos de representación: Sistema IRAM.
                \item Sistema de representación: perspectiva caballera, perspectiva isométrica 
                \item Rebatimientos y desarrollos.
                
            \end{itemize}
        \end{minipage} &  \begin{minipage} [t] {0.4\textwidth} 
            \begin{itemize}
                \item Representar un cuerpo por medio de vistas, en los distintos métodos de representación.
                \item Dadas las vistas principales de un cuerpo, identificar en qué método está representado.
                \item Dibujar un cuerpo en perspectiva dadas sus vistas principales.
            \end{itemize}
        \end{minipage} \\ \hline
       \end{tabular}
    \caption{Bloque 2}
\end{table}

\begin{table}[h]
    \centering
    \begin{tabular}{|c|c|} \hline
         \multicolumn{2}{|c|}{Contenidos} \\ \hline
         Conceptuales & Procedimentales  \\ \hline \hline
        \begin{minipage} [t] {0.4\textwidth} 
            \begin{itemize}
                \item Acotaciones.
                \item Acotación de los dibujos. 
                \item Clasificación de las cotas.
                \item Ejecución de líneas de cota y de referencia.
                \item Sistemas de acotamiento.
                \item Acotaciones condicionales a líneas base o de referencia
                \item Dimensiones que no es preciso acotar.
                \item Dimensiones de piezas normalizadas.
                \item Modos de acotar.
                \item Aplicaciones diversas.
                \item Símbolos.
            \end{itemize}
        \end{minipage} &  \begin{minipage} [t] {0.4\textwidth} 
            \begin{itemize}
                \item Acotado de un dibujo.
                \item Representación de dibujos en láminas acotados y en escala correspondiente.
                \item Confección de croquis.
            \end{itemize}
        \end{minipage} \\ \hline
       \end{tabular}
    \caption{Bloque 3}
\end{table}
\begin{table}[h!]
    \centering
    \begin{tabular}{|c|c|} \hline
         \multicolumn{2}{|c|}{Contenidos} \\ \hline
         Conceptuales & Procedimentales  \\ \hline \hline
        \begin{minipage} [t] {0.4\textwidth} 
            \begin{itemize}
                \item Generalidades del software FreeCAD y AutoCAD.
                \item Estructura de pantalla.
                \item Vistas y manejo de grilla.
                \item Modos de seleccionar objetos.
                \item Comandos:línea, rectángulos, polígonos,copiar, pegar, recortar,
                rotar,mover,espejar,trazar paralelas,empalmar líneas, extender,borrar, escribir un texto,distancia,área, círculos y arcos, zoom.
                \item Editar un texto.
                \item Cargar distintos tipos de línea.
                \item Generar extrusiones.
                \item Generar Objetos en 3D.
                \item Acotar e imprimir.
                \item Impresión de trabajos.
                
            \end{itemize}
        \end{minipage} &  \begin{minipage} [t] {0.4\textwidth} 
            \begin{itemize}
                \item Dibujo de formato de hoja y rótulo, en software freeCAD.
                \item Ejecución en FreeCAD, de los dibujos de la guía de trabajos prácticos.
                \item Impresión de planos realizados.
                \item Aplicación de conocimientos en el resto de los bloques.
                \item Uso de plotter e impresora.
            \end{itemize}
        \end{minipage} \\ \hline
       \end{tabular}
    \caption{Bloque 4}
\end{table}

\begin{table}[h!]
    \centering
    \begin{tabular}{|l|p{5cm}|p{5cm}|} \hline
        TÉCNICAS & INSTRUMENTOS & SITUACIONES DE EVALUACIÓN \\ \hline \hline
        OBSERVACIÓN & Lista de control & Realización de tareas\\ \hline
        ANÁLISIS DE TAREAS & Carpeta & Actividades no regladas\\ \hline
        OBSERVACIÓN & Registro de conductas grupales, informes, trabajos & Trabajo en Equipo, puesta en común\\ \hline 
        PRUEBAS & Prueba de proceso escrita individual & Participación en el diseño de investigación\\ \hline
    \end{tabular}
    \caption{Técnicas e instrumentos de evaluación}
    
\end{table}
\end{document}
